% =====================================================================
%  APPENDIX B — ASI measurement: methodology, formulas and assumptions
%  Inserted between the ERP appendix and the supplementary figures by
%  build_downstream_report.py. Referenced from Section 4 as Appendix~\ref{app:asi}.
% =====================================================================

\section{Appendix: ASI Measurement --- Segmentation Rule, Formulas and Assumptions}
\label{app:asi}

This appendix sets out in full the segmentation rule, estimators and assumptions
underlying Section~\ref{sec:asi_method} and the figures that follow it. It is
placed here rather than in the body so that the argument can be read without
interruption; nothing in Section~4 depends on material that is not stated here.

\subsection{The segmentation rule}
\label{app:asi_rule}

Following Corden \cite{corden1966}, the unit of effective-protection analysis is an
\textbf{activity}: a process converting traded inputs into a traded output, whose
protection depends jointly on the tariff on that output and the tariffs on its
inputs. Formally, for an activity $j$ with output tariff $t_j$ and input tariffs
$t_i$ over an input set with cost shares $a_{ij}$ (valued at free-trade prices),
the effective rate of protection is
\begin{equation}
    \mathrm{ERP}_j \;=\; \frac{t_j - \sum_i a_{ij} t_i}{1 - \sum_i a_{ij}} .
    \label{eq:erp}
\end{equation}
Equation~\eqref{eq:erp} makes two properties explicit that the body of the report
relies on. First, $\mathrm{ERP}_j < 0$ whenever $\sum_i a_{ij} t_i > t_j$: a
sufficiently protected input drives effective protection on the output negative,
which is the inverted-duty-structure case. Second,
$\partial \mathrm{ERP}_j / \partial t_i = -a_{ij}/(1 - \sum_i a_{ij})$, so the
damage done by an input tariff rises with the input-cost share $\sum_i a_{ij}$ ---
equivalently, it rises as the share of value added in output falls. This is the
formal content of the ``transmission mechanism'' paragraph in Section~4.7:
downstream fabrication carries an input-cost share of 83.1 per cent, so the
denominator of \eqref{eq:erp} is small and any input-tariff wedge is amplified.
The operational estimator actually used on Supply--Use data is the modified
Corden form given in Appendix~\ref{app:erp}.

The classification rule follows directly. Let $\mathcal{H}$ denote the set of HS
headings on which Section~\ref{sec:erp} estimates effective protection, that is
$\mathcal{H} = \{7604,\dots,7616\}$, and let $h(i)$ be the principal output
heading of NIC class $i$. Then
\begin{equation}
    i \in D \iff h(i) \in \mathcal{H},
    \qquad
    i \in U \iff h(i) \in \{7601, 7602, 7603\},
    \label{eq:rule}
\end{equation}
giving $U = \{2420\}$ and $D = \{2432,\, 2511,\, 2599\}$. Classes whose principal
output lies outside HS Chapter~76 --- NIC~2710, 2732, 2733 and 3030, whose outputs
fall in HS Chapters~85 and~88 --- are excluded: they are distinct activities in the
sense of \eqref{eq:erp}, facing a different output tariff and a different input
vector. The concordance produced by \eqref{eq:rule} is Table~\ref{tab:nic_concordance}.

\subsection{Notation}

For NIC class $i$ and financial year $t$, the ASI reports
\[
\begin{array}{ll}
O_{it} & \text{total output} \\
M_{it} & \text{total inputs} \\
W_{it} & \text{wages and salaries} \\
L_{it} & \text{total persons engaged}
\end{array}
\qquad
\begin{array}{ll}
N_{it} & \text{total workers} \\
C_{it} & \text{workers employed through contractors} \\
\pi_{it} & \text{net profit} \\
\tau_{t} & \text{effective corporate tax rate in year } t
\end{array}
\]
For a segment $S \in \{U, D\}$ all quantities aggregate additively,
$X_{St} = \sum_{i \in S} X_{it}$.

\subsection{Estimators}

Gross value added is the published accounting identity
\begin{equation}
    V_{it} \;=\; O_{it} - M_{it},
    \label{eq:gva}
\end{equation}
which holds in the source data to within $10^{-6}$ per cent in every class and
year. Value added per person engaged, and the labour-intensity ratio on which
Section~4.3 rests, are
\begin{equation}
    v_{St} \;=\; \frac{V_{St}}{L_{St}},
    \qquad
    \lambda_t \;=\; \frac{v_{Ut}}{v_{Dt}} .
    \label{eq:lambda}
\end{equation}
$\lambda_t > 1$ means a rupee of value added created downstream sustains
$\lambda_t$ times the employment of a rupee created upstream. The remaining
shares are
\begin{equation}
    \omega_{St} = \frac{W_{St}}{V_{St}}
    \quad\text{(labour share)},
    \qquad
    c_{St} = \frac{C_{St}}{N_{St}}
    \quad\text{(contract share)},
    \qquad
    m_{St} = \frac{M_{St}}{O_{St}}
    \quad\text{(input-cost share)} .
    \label{eq:shares}
\end{equation}

\paragraph{Corporate tax and employment yield.} Corporate tax is not published by
the ASI and is estimated as net profit multiplied by the effective corporate tax
rate for the year, with loss-making observations excluded rather than netted off:
\begin{equation}
    T_{St} \;=\; \tau_t \sum_{i \in S} \pi_{it}\,\mathbf{1}\{\pi_{it} > 0\},
    \qquad
    y_{St} \;=\; \frac{L_{St}}{T_{St}} ,
    \label{eq:tax}
\end{equation}
where $y_{St}$ is persons engaged per unit of corporate tax. The indicator
$\mathbf{1}\{\cdot\}$ matters: a negative $\pi$ would otherwise generate a negative
tax, which has no behavioural interpretation. One observation within the retained
classes is affected (NIC~2432 in 2013--14).

\paragraph{Index, growth and volatility.} With base year $b = 2011\text{--}12$,
\begin{equation}
    I_{St} \;=\; 100 \times \frac{V_{St}}{V_{Sb}} .
    \label{eq:index}
\end{equation}
Because $V_{Ut}$ is volatile, growth is reported on three-year endpoint averages
as well as on single endpoints. Writing $A = \{2011\text{--}12,\,2012\text{--}13,\,2013\text{--}14\}$
and $B = \{2020\text{--}21,\,2021\text{--}22,\,2022\text{--}23\}$, and $\bar V_{S\mathcal{A}}$
for the mean of $V_{St}$ over a set $\mathcal{A}$,
\begin{equation}
    g_S \;=\; \frac{\bar V_{SB}}{\bar V_{SA}} - 1,
    \qquad
    \mathrm{CAGR}_S \;=\; \left( \frac{\bar V_{SB}}{\bar V_{SA}} \right)^{1/9} - 1 ,
    \label{eq:growth}
\end{equation}
the exponent being the nine years separating the midpoints of $A$ and $B$.
Volatility is the sample standard deviation of year-on-year log growth over the
twelve clean years,
\begin{equation}
    \sigma_S \;=\; \mathrm{sd}\!\left( \ln V_{S,t+1} - \ln V_{St} \right) .
    \label{eq:vol}
\end{equation}

\subsection{Construction of each figure}
\label{app:asi_figures}

Figure notes in the body carry no derivations. The construction of each is given
here. Throughout, $S \in \{U, D\}$, $t$ indexes the financial year, and the base
year is $b = 2011\text{--}12$. Series are computed on the twelve-year window
$t \in [2011\text{--}12,\ 2022\text{--}23]$ except where stated.

\begin{table}[H]
    \centering
    \small
    \begin{tabular}{llp{7.8cm}}
        \toprule
        \textbf{Figure} & \textbf{Estimator} & \textbf{Construction and domain} \\
        \midrule
        4.1 & \eqref{eq:index} & $I_{St} = 100\,V_{St}/V_{Sb}$, both segments, $t$ extended
        to 2023--24. Current prices; no deflator. \\
        \addlinespace
        4.2 & \eqref{eq:lambda} & $\lambda_t = v_{Ut}/v_{Dt}$. Reference lines: the sample
        median of $\{\lambda_t\}$, and the exogenous 2.5$\times$ benchmark of
        \cite{mom_vision2025}, which is not an estimate of ours. \\
        \addlinespace
        4.3 & \eqref{eq:shares} & $\omega_{St} = W_{St}/V_{St}$. Numerator and denominator
        are both monetary, so the series is invariant to the treatment of persons engaged. \\
        \addlinespace
        4.4 & \eqref{eq:index} & $I_{it} = 100\,V_{it}/V_{ib}$ at class rather than segment
        level, $i \in \{2420, 2432, 2511, 2599\}$. The horizontal reference is $I = 100$. \\
        \addlinespace
        4.5 & \eqref{eq:shares} & $c_{St} = C_{St}/N_{St}$. The denominator is total
        \emph{workers} $N$, not persons engaged $L$: the direct/contract partition is
        defined only within $N$. \\
        \addlinespace
        4.6 & \eqref{eq:tax} & $y_{St} = L_{St}/T_{St}$. Since $\tau_t$ is common to both
        segments it cancels from $y_{Dt}/y_{Ut}$, so the ratio is invariant to the
        assumed tax rate. \\
        \addlinespace
        4.7 & --- & Four ratios evaluated at $t = 2022\text{--}23$, each formed from
        estimators already defined: $v_U/v_D = 24.97/7.95$; $\omega_D/\omega_U =
        38.41/17.59$; $y_D/y_U = 150.83/38.73$; $\sigma_U/\sigma_D = 0.3732/0.1046$.
        The fourth is oriented upstream-over-downstream because it is upstream value
        added that is the more variable. \\
        \bottomrule
    \end{tabular}
\end{table}

\noindent \textbf{Domain restrictions.} All series terminate at 2022--23 under
assumption~A7 below. Figure~4.1 extends to 2023--24 because the index depends only on
$V$, which is unaffected by the two breaks in that release; every series involving
$W$ or the direct/contract partition necessarily stops at 2022--23. Loss-making
observations are excluded from $T_{St}$ by the indicator in \eqref{eq:tax}; one
observation within the retained classes is affected.

\subsection{Two derived results used in Section~4}

\paragraph{Amplification of an input-cost wedge.} Let $m = M/O$ be bought-in inputs as
a share of output value and $\theta$ the share of the input bill accounted for by
metal. Since $V = (1-m)O$, a proportional rise $\hat p$ in the metal price reduces
value added by
\begin{equation}
    \frac{\Delta V}{V} \;=\; -\,\frac{m\theta}{1-m}\,\hat p ,
    \label{eq:passthrough}
\end{equation}
so that $m/(1-m)$ is the factor by which a proportional input-cost change is
amplified when it reaches value added. This is equation~\eqref{eq:erp} viewed from
the producer's side: the denominator $1 - \sum_i a_{ij}$ that governs the effective
rate is the same value-added share that governs the amplification. With the measured
downstream input share $m = 0.8309$ the factor is \textbf{4.91}; upstream, where the
metal is produced rather than purchased, $\theta \approx 0$ and the effect vanishes.

\begin{table}[H]
    \centering
    \small
    \begin{tabular}{ccc}
        \toprule
        \textbf{Metal share of} & \textbf{Value added lost} & \textbf{Value added lost} \\
        \textbf{input bill, $\theta$} & \textbf{per 1\% on metal} & \textbf{at a 7.5\% wedge} \\
        \midrule
        30\% & 1.48\% & 11.1\% \\
        40\% & 1.97\% & 14.7\% \\
        50\% & 2.46\% & 18.4\% \\
        60\% & 2.95\% & 22.1\% \\
        70\% & 3.44\% & 25.8\% \\
        \bottomrule
    \end{tabular}

    \vspace{0.15cm}
    \raggedright
    \small $\theta$ is not observable in the ASI, which records the input bill in
    value but not by material. The range spans plausible values for aluminium
    conversion; firm-level cost data would replace it with a point estimate.
\end{table}

\paragraph{Employment per tonne.} Writing $Q$ for tonnes processed, the per-tonne
employment ratio decomposes exactly into the per-rupee ratio and the ratio of value
added per tonne:
\begin{equation}
    \frac{L_D/Q_D}{L_U/Q_U}
    \;=\;
    \underbrace{\frac{L_D/V_D}{L_U/V_U}}_{\lambda \;=\; 3.14}
    \times
    \underbrace{\frac{V_D/Q_D}{V_U/Q_U}}_{1.81}
    \;=\; 5.7 .
    \label{eq:pertonne}
\end{equation}
Both factors exceed unity, so the two normalisations reinforce rather than offset one
another. Tonnage is not an ASI variable: $Q_U$ is installed primary capacity for FY24
and $Q_D$ an industry estimate of 2.0--2.5 Mt, which is why Section~4.3 reports the
per-tonne figure as a range of 5.1--6.4 and treats it as supporting rather than
primary evidence. Equation~\eqref{eq:pertonne} also disposes of the much larger
ratios sometimes quoted: a figure near 35:1 requires counting the direct workforce of
the primary smelters alone against the direct, indirect and unorganised workforce
downstream, and the two mismatches together account for the whole of the difference.

\subsection{Apportionment, and why it does not affect the results}
\label{app:asi_alpha}

None of the classes in $U$ or $D$ is exclusive to aluminium. The natural objection
is that each should be scaled by an aluminium share $\alpha_i \in (0,1]$, giving
$V^{Al}_{i} = \alpha_i V_{i}$ and $L^{Al}_{i} = \alpha_i L_{i}$. Applying a
coefficient derived from output shares to an employment variable embeds the
\emph{industry technology assumption} --- that aluminium activity within a class has
the same factor intensity as the rest of that class --- as against the
\emph{commodity technology assumption}, under which each product carries its own
factor requirements wherever it is made \cite{millerblair}. The choice is
consequential in general. Here it is not, for a reason that can be stated exactly
rather than simulated.

\paragraph{Proposition.} \textit{For any apportionment vector
$\alpha = (\alpha_i)_{i \in D}$ with $\alpha_i > 0$, the apportioned segment
productivity $v_D(\alpha)$ is a weighted average of the class productivities
$v_i = V_i / L_i$, and is therefore bounded by their range irrespective of
$\alpha$.}

\paragraph{Proof.} Writing $w_i(\alpha) = \alpha_i L_i \big/ \sum_{j \in D} \alpha_j L_j$,
\begin{equation}
    v_D(\alpha)
    \;=\; \frac{\sum_{i \in D} \alpha_i V_i}{\sum_{i \in D} \alpha_i L_i}
    \;=\; \sum_{i \in D} \frac{\alpha_i L_i}{\sum_{j} \alpha_j L_j} \cdot \frac{V_i}{L_i}
    \;=\; \sum_{i \in D} w_i(\alpha)\, v_i ,
    \label{eq:weighted}
\end{equation}
with $w_i(\alpha) > 0$ and $\sum_i w_i(\alpha) = 1$. Hence
$\min_i v_i \le v_D(\alpha) \le \max_i v_i$. Since $U$ is a single class,
$v_U$ is exactly invariant to $\alpha_U$, which cancels in \eqref{eq:lambda}.
Therefore
\begin{equation}
    \frac{v_U}{\max_{i \in D} v_i} \;\le\; \lambda(\alpha) \;\le\; \frac{v_U}{\min_{i \in D} v_i}
    \qquad \text{for every } \alpha . \qquad \rule{0.9ex}{0.9ex}
    \label{eq:bound}
\end{equation}

\paragraph{The bound, evaluated.} In 2022--23 the class productivities are
$v_{2432} = 7.43$, $v_{2511} = 7.38$ and $v_{2599} = 8.30$ lakh per person, against
$v_U = 24.97$. Substituting into \eqref{eq:bound},
\[
    3.01 \;\le\; \lambda(\alpha) \;\le\; 3.38
    \qquad \text{for \emph{any} apportionment whatever,}
\]
against the reported undiscounted value of 3.14. The result requires no assumption
about the magnitude of any $\alpha_i$ --- only that the shares be positive. It also
explains why the 48-scenario grid reported in Section~4.9 moves $\lambda$ so little:
every one of those scenarios necessarily lies inside \eqref{eq:bound}.

\paragraph{What is \emph{not} invariant.} The same algebra shows that
\emph{levels} carry no such protection. The cross-segment employment multiple
$L_D(\alpha)/L_U(\alpha) = \sum_i \alpha_i L_i / (\alpha_U L_U)$ is homogeneous of
degree one in $\alpha$ and ranges from 0.81 to 3.07 across plausible vectors.
This replicates, in a different setting, Bohlin and Widell's finding that the
technology assumption matters greatly for technical coefficients and little for
factor-content applications \cite{bohlin2006}. Accordingly the report quotes
intensity ratios, shares and per-rupee statistics throughout, and makes no
cross-segment comparison of employment levels.

\subsection{Assumptions}
\label{app:asi_assumptions}

The following are stated so that each can be checked or contested independently.

\begin{enumerate}[label=\textbf{A\arabic*.}, leftmargin=2.6em]
    \item \textbf{Partition criterion.} Segments are defined by \eqref{eq:rule}: a
    class is downstream if and only if its principal output lies in HS~7604--7616.
    This is Corden's activity criterion, not a process-stage criterion; the
    material-flow literature separates semi-fabrication from fabrication, and the
    consequence of adopting that alternative is reported in Section~4.9.

    \item \textbf{Nominal valuation.} All values are at current prices with no
    deflator applied. Growth rates are therefore nominal and overstate real
    expansion; comparisons between segments are unaffected, since both are
    measured on the same basis.

    \item \textbf{Principal-product classification.} The ASI assigns an
    establishment wholly to the NIC class of its principal product. Because
    NIC~2420's explanatory note includes semi-manufacturing, foil and wire
    drawing, output of HS~7604--7608 produced within integrated smelters is
    recorded upstream. This biases every comparison in the report \emph{against}
    its own conclusion; no correction is applied.

    \item \textbf{No apportionment by metal.} Class aggregates are used
    undiscounted, so segment totals measure the registered manufacturing space
    within which aluminium activity sits, not a pure aluminium account.
    Appendix~\ref{app:asi_alpha} shows the reported ratios are bounded
    irrespective of any apportionment.

    \item \textbf{Uniform tax rate.} The effective corporate tax rate $\tau_t$ in
    \eqref{eq:tax} is applied uniformly across classes within a year. Differences
    in effective rates between segments are not observable in the published data.

    \item \textbf{Loss exclusion.} Loss-making observations are excluded from the
    tax series rather than netted off, per \eqref{eq:tax}.

    \item \textbf{Terminal year.} 2023--24 is excluded from all stated statistics.
    Reported wages fall in every class in that year while output and persons
    engaged rise, and the identity $N_{it} = $ directly employed $+$ contract fails
    with a residual of 178,591 workers (19.6 per cent of the total), having held
    exactly in each of the twelve preceding years.

    \item \textbf{Frame coverage.} The ASI covers registered factories only ---
    units with 10 or more workers using power, or 20 or more without. Downstream
    aluminium fabrication has a substantial unincorporated segment outside this
    frame; upstream smelting has none. The downstream aggregate is therefore an
    undercount to an extent the published data cannot quantify.

    \item \textbf{Missing classes.} NIC~2512 (tanks, reservoirs, gas containers) is
    not disaggregated in the ASI All-India summary, so HS~7611 and 7613 --- two of
    the thirteen ERP lines --- have no ASI counterpart and are absent from the
    downstream aggregate.
\end{enumerate}

Assumptions A3, A8 and A9 all run in the same direction: each inflates the
measured upstream segment or deflates the measured downstream one. The gaps
reported in Section~4 are accordingly lower bounds.

\subsection{Robustness}
\label{app:asi_robust_table}

Every quantitative claim in Section~4 was tested against the objections a reader is
likely to raise. The tests are set out here rather than in the body because several of
them are only interesting if the answer is unexpected, and none of them is. Two are
findings in their own right and are flagged as such.

\paragraph{Sign consistency: nothing depends on the year chosen.} Across the twelve
clean years, and without a single exception, downstream fabrication is the more
labour-intensive segment ($\lambda_t > 1$); the ratio exceeds the Ministry of Mines'
2.5$\times$ lower bound; the downstream wage share exceeds the upstream wage share
($\omega_{Dt} > \omega_{Ut}$); downstream persons engaged exceed upstream
($L_{Dt} > L_{Ut}$); and downstream employment per unit of estimated tax exceeds
upstream ($y_{Dt} > y_{Ut}$). Twelve out of twelve on all five. This is the strongest
single robustness result in the section: no headline claim is an artefact of the year
selected.

\paragraph{Terminal year.} Recomputing $\lambda$ to alternative terminal years gives
3.17, 3.08, 3.75, 5.81, 3.14 and 3.05 for 2018--19 through 2023--24. The value
reported in the body, 3.14, sits near the bottom of that range --- which is the answer
to any suggestion that the terminal year was chosen to flatter.

\paragraph{Base year.} Single-endpoint growth comparisons \emph{are} sensitive to the
base year, because $V_{Ut}$ is volatile; the sign of the growth differential can be
moved by the choice. This is why growth is also reported on the three-year endpoint
averages of \eqref{eq:growth}, which are stable: upstream 105.6 per cent against
downstream 73.0 per cent, compound rates of 8.34 and 6.28 per cent. Both bases give
the same ordering and a similar gap, and no claim in Section~4 rests on the
single-endpoint comparison alone.

\paragraph{Leave-one-out.} Dropping each downstream class in turn moves the 2022--23
labour-intensity ratio to 3.13 (without 2432), 3.04 (without 2511) and 3.38 (without
2599), against 3.14 for the full basket. No single class drives the result.

\paragraph{The process-taxonomy objection.} Section~\ref{app:asi_rule} groups
semi-fabrication with fabrication on Corden's criterion, whereas the material-flow
literature separates them as distinct process steps \cite{cullen2013,bertram2009}.
Because NIC~2432 is the semis-and-casting class, the objection can be settled by
reassigning it and recomputing. Doing so moves $\lambda$ from 3.14 to 2.77, the
wage-share gap from 20.8 to 19.6 percentage points, and the employment-yield ratio
from 3.89 to 3.31. Every conclusion survives; only the magnitudes move, and modestly.
The reason is visible in the underlying figure: NIC~2432's own value added per person
engaged is INR 7.43 lakh, close to the INR 7.99 lakh of the remaining fabrication
classes and nowhere near the INR 24.97 lakh of primary metal. On the measure that
matters here, semi-fabrication behaves like fabrication rather than like smelting,
whichever taxonomy one prefers --- which is a more durable position than winning the
classification argument would have been.

\paragraph{Basket width.} Earlier drafts used a wider downstream basket, adding
NIC~2710 (motors, generators, transformers) and 2732 (insulated wire and cable).
Those classes fail the rule of \eqref{eq:rule} --- their outputs are HS Chapter~85, a
different activity protected by different tariffs --- but the findings do not depend on
excluding them. On the wide basket $\lambda$ is 2.61 rather than 3.14, the wage-share
gap 18.9 rather than 20.8 percentage points, and the employment-yield ratio 2.91
rather than 3.89. Every direction is unchanged; the magnitudes move \emph{against} the
wider basket. The narrow basket is preferred not because it is more favourable but
because it is the only one consistent with the HS~7604--7616 scope on which the rest of
the report is built.

\paragraph{Summary.} Table~\ref{tab:asi_robust} collects the results. The reported
specification is the first row of each block.

\begin{table}[H]
    \centering
    \caption{Sensitivity of the labour-intensity ratio $\lambda$ and associated statistics}
    \label{tab:asi_robust}
    \resizebox{\textwidth}{!}{
    \begin{tabular}{lrrr}
        \toprule
        \textbf{Specification} & $\boldsymbol{\lambda}$ & \textbf{Wage-share gap} & \textbf{Jobs per rupee of tax} \\
        \midrule
        \multicolumn{4}{l}{\textit{Reported specification}} \\
        Upstream 2420; downstream 2432, 2511, 2599 & \textbf{3.14} & \textbf{20.8 pp} & \textbf{3.89} \\
        \midrule
        \multicolumn{4}{l}{\textit{Terminal year}} \\
        2018--19 & 3.17 & --- & --- \\
        2019--20 & 3.08 & --- & --- \\
        2020--21 & 3.75 & --- & --- \\
        2021--22 (metal-price peak) & 5.81 & --- & --- \\
        2023--24 (excluded, shown for reference) & 3.05 & --- & --- \\
        \midrule
        \multicolumn{4}{l}{\textit{Leave-one-out on the downstream basket}} \\
        Excluding 2432 & 3.13 & --- & --- \\
        Excluding 2511 & 3.04 & --- & --- \\
        Excluding 2599 & 3.38 & --- & --- \\
        \midrule
        \multicolumn{4}{l}{\textit{Alternative segment definitions}} \\
        Material-flow reading: 2432 grouped upstream & 2.77 & 19.6 pp & 3.31 \\
        Wide basket: adding 2710 and 2732 & 2.61 & 18.9 pp & 2.91 \\
        \midrule
        \multicolumn{4}{l}{\textit{Apportionment}} \\
        Analytic bound over all $\alpha$, eq.~\eqref{eq:bound} & 3.01--3.38 & --- & --- \\
        48-scenario grid & 3.19--3.21 & --- & --- \\
        \bottomrule
    \end{tabular}
    }

    \vspace{0.2cm}
    \raggedright
    \small Source: our calculations from ASI All-India summary results, 2022--23 unless
    otherwise stated. ``---'' denotes a statistic not recomputed for that variant because it
    is not at issue. Every row preserves the sign and the order of magnitude of the reported
    specification; the lowest value anywhere in the table, 2.61, still exceeds the Ministry of
    Mines' 2.5$\times$ lower bound.
\end{table}


